\BeleskeID{01-s037}
% element: 01-s037:e01
% element: 01-s037:e02
% element: 01-s037:e03
% element: 01-s037:e04

% element: 01-s037:d01
Označimo proizvode prve dve grane sa \(a=x_1x_2\) i \(b=x_3x_4\). Potrebna funkcija je \(a+b\). Dualni prikaz poslednjeg kola računa isti zbir kao negirani proizvod komplementarnih ulaza:
\[
a+b=\overline{\bar{a}\,\bar{b}}.
\]
Ulazna NI kola neposredno stvaraju \(\bar{a}\) i \(\bar{b}\). Poslednje NI zato vraća traženi zbir proizvoda, bez zasebnih invertora na međuvezi. Negacioni kružići u drugom crtežu opisuju funkciju poslednjeg elementa, dok u trećem njegove potrebne ulazne vrednosti već obezbeđuju prethodna kola.

Korišćenje dva različita tipa kola može biti praktično nezgodno i skuplje. Tip pri izboru obuhvata i funkciju i broj ulaza, pa isti naziv operacije ne znači automatski isti tip komponente. Realizacija samo dvoulaznim NI kolima omogućava upotrebu jednog tipa; izbor ipak zavisi od konkretnih raspoloživih komponenti.
